\section{Foundations}\label{foundations}

\subsection{Two Credibility Domains}\label{two-credibility-domains}

This paper distinguishes two fundamentally different credibility domains that obey different dynamics:

\textbf{Definition 2.0a (Mathematical Credibility).} \emph{Mathematical credibility} $C_M$ measures the probability that a claim is logically sound. The audience is a formal verifier (proof assistant, compiler, test suite). The signal space consists of artifacts that can be mechanically checked. Mathematical credibility is binary at the limit: a proof compiles or it doesn't.

\textbf{Definition 2.0b (Social Credibility).} \emph{Social credibility} $C_S$ measures the probability that a social audience accepts a claim. The audience is human agents with priors shaped by institutional hierarchy, reputation, and group membership. The signal space consists of credentials, affiliations, endorsements, and communication patterns.

\textbf{Theorem 2.0c (Domain Independence).} Mathematical credibility and social credibility are orthogonal:
\[C_M(c, s) \not\Rightarrow C_S(c, s) \quad \text{and} \quad C_S(c, s) \not\Rightarrow C_M(c, s)\]
A signal maximizing one domain does not necessarily affect the other.

\emph{Proof.} Constructive. (1) High $C_M$, low $C_S$: a correct proof by an unknown author receives $C_M \to 1$ (it compiles) but may receive $C_S \approx P_{\text{prior}}$ (no institutional endorsement). (2) High $C_S$, low $C_M$: a claim by a prestigious institution without formal verification receives high $C_S$ (reputation transfer) but $C_M$ is undefined or low (no proof). \qed

\textbf{Corollary 2.0d (Costly Signal Domain-Specificity).} Costly signals are domain-specific:
\begin{itemize}
\item Machine-checked proofs are maximally costly in the mathematical domain (cannot compile a false proof) but may be cheap talk in the social domain (anyone can claim to have proofs).
\item Institutional credentials (PhD, tenure, affiliation) are costly in the social domain (years of compliance) but cheap talk in the mathematical domain (credential $\not\Rightarrow$ soundness).
\end{itemize}

\textbf{Remark (Credibility Domain Conflict).} When a signal achieves $C_M \to 1$ but $C_S \approx 0$, the mathematical and social domains are in conflict. A rational response in the mathematical domain (engage with proofs) differs from a rational response in the social domain (defer to hierarchy). Observers may respond in either domain. This paper's theorems apply within each domain separately; cross-domain dynamics require modeling both simultaneously.

\subsection{Signals and Costs}\label{signals-and-costs}

\textbf{Definition 2.1 (Signal).} A \emph{signal} is a tuple
\(s = (c, v, p)\) where: \(c\) is the \emph{content} (what is
communicated), \(v \in \{\top, \bot\}\) is the \emph{truth value}
(whether content is true), and \(p \in \mathbb{R}_{\geq 0}\) is the
\emph{production cost}.

\textbf{Physical grounding of cost.}
Production cost is measured in joules. By Landauer's principle \cite{landauer1961}, each
irreversible bit operation erases at minimum
\[
E_{\min} = k_B T \ln 2 \;\approx\; 2.8 \times 10^{-21}~\text{J at } T = 300\,\text{K},
\]
where $k_B$ is Boltzmann's constant. Paper 4's Energy Lower Bound (Theorem 13) shows that
srank bits of structural information carry a minimum physical energy cost
proportional to $\text{srank} \cdot k_B T \ln 2$. Signal production cost inherits this floor.

\textbf{Definition 2.2 (Cheap Talk).} A signal \(s\) is \emph{cheap
talk} if production cost is truth-independent:
\[\text{Cost}(s | v = \top) = \text{Cost}(s | v = \bot)\]
In physical units: the joules expended to produce $s$ do not depend on whether $v = \top$ or $v = \bot$.
Typed text is cheap talk: the Landauer cost of typing ``true'' equals the cost of typing ``false.''

\textbf{Definition 2.3 (Costly Signal).} A signal \(s\) is \emph{costly}
if:
\[\Delta := \text{Cost}(s | v = \bot) - \text{Cost}(s | v = \top) > 0\]
Producing the signal when false costs strictly more than when true. The quantity $\Delta$
(in joules) is the \emph{cost differential}.

\textbf{Definition 2.3a (Credibility as Perceived Integrity via Heuristic).}
A \emph{credibility heuristic} is a polynomial-time computable function 
$h : \mathcal{C} \times \mathcal{S}^* \to [0,1]$ that maps a claim and observed 
signals to a perceived integrity score. The \emph{credibility} of a claim $c$ given 
signals $s$ is defined as:
\[C(c, s) := h(c, s)\]
where $h$ is the receiver's heuristic for assessing claim integrity.

\textbf{Rationale.} This definition makes credibility operational: it is not an 
abstract notion but a \emph{perceived} assessment computed by a specific \emph{heuristic}. 
The receiver's cognitive system uses shortcuts to evaluate whether a claim's 
proponent is trustworthy, competent, or sincere. This grounds the subsequent 
analysis in algorithmic information processing rather than undefined intuition.

\textbf{Heuristic Properties.} We assume $h$ satisfies: (i) $h(c, \emptyset)$ returns 
the receiver's prior (baseline integrity belief), (ii) $h$ is efficient (bounded 
computation), and (iii) $h$ may exhibit systematic bias (deviation from Bayesian optimality).

\textbf{Mathematical vs Social Domain.} The heuristic differs by domain:
\begin{itemize}
\item $h_M$: mathematical credibility heuristic (type-checker, proof verifier) 
      - converges to 0/1 as verification runs (bounded verification $\to$ bounded error)
\item $h_S$: social credibility heuristic (reputation, credentials, rhetorical style)
      - bounded away from extremes due to cognitive limitations (bounded memory $\to$ bounded precision)
\end{itemize}

\textbf{Remark (Lean proofs as physical signals).}
A machine-checked Lean proof is a costly signal with $\Delta \to \infty$. To produce a
compiling proof of a false theorem, the type-checker's false-positive rate satisfies
$\varepsilon_F \to 0$ (soundness). By Paper 4's Counting Gap (Theorem 3):
$\varepsilon \cdot N \leq C \Rightarrow N \leq C / \varepsilon$. With $\varepsilon_F \to 0$,
the number of false proofs that pass verification is bounded by $C \cdot \varepsilon_F \to 0$.
Producing one would require traversing a search space of diverging size, hence diverging energy.
Therefore $\Delta_{\text{Lean}} = \infty$: the Lean type-checker physically separates true from false.

\textbf{Intuition:} Verbal assertions are cheap talk---saying ``I'm
honest'' costs the same whether you're honest or not. A PhD from MIT is
a costly signal \cite{spence1973job}---obtaining it while incompetent is much harder than
while competent. Similarly, price and advertising can serve as signals of quality \cite{milgrom1986price}.

\subsection{Credibility Functions}\label{credibility-functions}

\textbf{Definition 2.4 (Prior).} A \emph{prior} is a probability
distribution \(P : \mathcal{C} \to [0,1]\) over claims, representing
beliefs before observing signals.

\textbf{Definition 2.5 (Credibility Function).} A \emph{credibility
function} is a mapping:
\[C : \mathcal{C} \times \mathcal{S}^* \to [0,1]\] from (claim,
signal-sequence) pairs to credibility scores, satisfying: 1.
\(C(c, \emptyset) = P(c)\) (base case: prior) 2. Bayesian update:
\(C(c, s_{1..n}) = P(c | s_{1..n})\)

\textbf{Definition 2.6 (Rational Agent).} An agent is \emph{rational}
if: 1. Updates beliefs via Bayes' rule 2. Has common knowledge of
rationality \cite{aumann1995backward} (knows others are rational, knows others know, etc.) 3.
Accounts for strategic signal production \cite{cho1987signaling}.

\subsection{Deception Model}\label{deception-model}

\textbf{Definition 2.7 (Deception Prior).} Let \(\pi_d \in [0,1]\) be
the prior probability that a random agent will produce deceptive
signals. This is common knowledge.

\textbf{Definition 2.8 (Magnitude).} The \emph{magnitude} of a claim
\(c\) is: \[M(c) = -\log P(c)\] High-magnitude claims have low prior
probability. This is the standard self-information measure \cite{shannon1948}.

\subsection{Model Contract and Regime Tags}\label{sec:credibility-contract}

All theorem statements in this paper are typed by the following contract:
\begin{itemize}
\item \textbf{K1 (binary claim state):} claim truth is modeled as \(C \in \{0,1\}\).
\item \textbf{K2 (Bayesian receiver):} posterior credibility is computed by Bayes updates.
\item \textbf{K3 (signal channel declaration):} each theorem declares whether the signal channel is cheap talk or verifier-backed.
\item \textbf{K4 (audience domain declaration):} each theorem declares mathematical-domain (\(C_M\)) or social-domain (\(C_S\)) scope.
\end{itemize}

We use the following regime tags:
\begin{itemize}
\item \texttt{[CT]}: cheap-talk channel (truth-independent signal cost, $\Delta = 0$),
\item \texttt{[VS]}: verifier-backed signal channel ($\varepsilon_T, \varepsilon_F$ model, $\Delta \to \infty$),
\item \texttt{[M]}: mathematical credibility domain,
\item \texttt{[S]}: social credibility domain.
\end{itemize}

\begin{proposition}[Declared Regime Coverage for Credibility Claims]\label{prop:credibility-regime-coverage}
Each declared regime above has at least one theorem-level mechanized core in Lean:
\texttt{[CT]} via \texttt{cheap\_talk\_bound}, \texttt{magnitude\_penalty}, \texttt{emphasis\_penalty};
\texttt{[VS]} via \texttt{verified\_signal\_credibility}, \texttt{proof\_as\_ultimate\_signal};
\texttt{[M]/[S]} via \texttt{domain\_independence\_math\_not\_implies\_social} and \texttt{domain\_independence\_social\_not\_implies\_math}.
The physical connection between \texttt{[CT]} ($\Delta=0$) and \texttt{[VS]} ($\Delta\to\infty$)
is grounded in Landauer's principle and Paper 4's Counting Gap (Theorem 3).
\end{proposition}

\begin{proof}
Direct by inspection of the theorem declarations in
\texttt{Credibility/CheapTalk.lean}, \texttt{Credibility/CostlySignals.lean}, and
\texttt{Credibility/Basic.lean}.
\end{proof}

\begin{center}\rule{0.5\linewidth}{0.5pt}\end{center}
